% answers_short.tex —— 2000 年·广东卷（科类未标） · 简短答案速查
% 由 tools/make_short_answers.py 自 answers.tex 的 \gksolhead/\gkanswer 提取，
% 并与 2000/2000_数学真题_广东_答案要点.md 逐题比对后生成，勿手改（改源 MD 后重跑）。

\begin{gknotice}
\textbf{使用说明}：本节只给最终答案便于快速核对，完整推导见解析版。答案为编写者独立推导，\textbf{非官方标准答案}。
\end{gknotice}

\begin{gknotice}
选择题答案中的 U 表示待核，未给出猜测答案，详情见解析版。
\end{gknotice}

\gkpart{选择题}
\begin{center}\begin{tabular}{ccccccccccccc}
\toprule
题号 & 1 & 2 & 3 & 4 & 5 & 6 & 7 & 8 & 9 & 10 & 11 & 12 \\
\midrule
答案 & A & B & D & D & D & C & B & C & A & C & C & U \\
\bottomrule
\end{tabular}\end{center}

\gkpart{填空题}
\begin{itemize}[leftmargin=2.6em]
  \item[\textbf{13.}] $252$
  \item[\textbf{14.}] $(-3/\sqrt5,3/\sqrt5)$
  \item[\textbf{15.}] $a_n=1/n$
  \item[\textbf{16.}] ②、③
\end{itemize}

\gkpart{解答题}
\begin{itemize}[leftmargin=2.6em]
  \item[\textbf{17.}] (1) $x=\pi/3+2k\pi$，$k\in\mathbb Z$；(2) 先左移 $\pi/6$，再沿纵轴伸长为原来的2倍
  \item[\textbf{18.}] (1) $a_1=1$，公比 $q=2$；(2) $T_n=2^{n+1}-n-2$
  \item[\textbf{19.}] (1) U/PENDING；(2) U/PENDING
  \item[\textbf{20.}] (1) $0<a<1$ 时 $[0,2a/(1-a^2)]$；$a\ge1$ 时 $[0,+\infty)$；(2) $f$ 在 $[0,+\infty)$ 上严格递减
  \item[\textbf{21.}] (1) $p(t)=300-t$（$0\le t\le200$），$p(t)=2t-300$（$200\le t\le300$）；$Q(t)=(t-150)^2/200+100$；(2) 图中 $t=50$ 天时纯收益最大
  \item[\textbf{22.}] U/PENDING
\end{itemize}
